\newpage
\pagestyle{plain}

\begin{center}
    \Large
    \textbf{Abstract}
\end{center}

\vspace{0.5cm}

Photoelectrochemical (PEC) water splitting is one of the most promising pathways for
solar hydrogen production, but the experimental search for high-performance
photoelectrode materials is slow and resource-intensive.
This work presents a full-stack machine learning system that predicts three key
performance metrics of PEC cells — Photocurrent Density (mA/cm\textsuperscript{2}),
Solar-to-Hydrogen (STH) efficiency (\%), and Hydrogen Evolution Rate
($\mu$mol\,h\textsuperscript{-1}\,cm\textsuperscript{-2}) — directly from
material and experimental parameters.

A curated dataset of 676 entries was assembled from 529 published experimental
results, augmented with physics-derived estimates and literature-based synthetic
rows to address severe class imbalance in STH and H\textsubscript{2} targets.
Twenty physics-informed features were engineered from raw inputs, including
bandgap--pH interaction, overpotential proxy, ionic strength proxy, and
Arrhenius-like temperature--pH kinetics terms.

An ensemble of XGBoost and Random Forest regressors is trained per target, with
hyperparameter selection via randomised search and generalisation measured by
5-fold cross-validation.
Best cross-validated R\textsuperscript{2} scores are 0.454 for Photocurrent,
0.533 for STH, and 0.366 for H\textsubscript{2} — compared to 0.065, 0.278,
and $-$0.090 before dataset augmentation.

A REST API (Flask) and single-page web interface expose the models for
interactive prediction, returning per-target values, model agreement scores,
and a composite performance classification (LOW / MEDIUM / HIGH) calibrated to
the training distribution.

\vspace{0.5cm}
\noindent\textbf{Keywords:} photoelectrochemical water splitting, machine learning,
XGBoost, Random Forest, hydrogen production, solar energy, materials informatics,
feature engineering, ensemble learning.
